
% Dateiname: MoE_Glossar_QUBLE_261029-1.tex

\documentclass[11pt,a4paper]{article}

\usepackage[T1]{fontenc}
\usepackage[utf8]{inputenc}
\usepackage[ngerman]{babel}
\usepackage[a4paper,margin=25mm]{geometry}
\usepackage{amsmath,amssymb}
\usepackage{helvet}
\usepackage{sfmath}
\renewcommand{\familydefault}{\sfdefault}
\usepackage{microtype}
\usepackage[hidelinks]{hyperref}
\usepackage{lastpage}
\usepackage{fancyhdr}
\usepackage{parskip}

\pagestyle{fancy}
\fancyhf{}
\fancyfoot[C]{\thepage\ (\pageref{LastPage})}
\renewcommand{\headrulewidth}{0pt}
\renewcommand{\footrulewidth}{0pt}

\hypersetup{
  pdftitle={MoE-Glossar: QUBLE / Quantelung},
  pdfauthor={}
}

\begin{document}

\vspace*{-1.8cm}

\noindent
\begin{minipage}[t]{0.66\textwidth}
{\LARGE\bfseries QUBLE / Quantelung}
\end{minipage}
\hfill
\begin{minipage}[t]{0.30\textwidth}
\raggedleft
{\small MoE-Glossareintrag}\\
{\small Version 261029-1}
\end{minipage}

\vspace{1.4em}

\section*{Definition}

Ein QUBLE ist die erste stabile räumliche
Dichtestruktur, die aus der Quantelung
des Apeirons hervorgeht.

Quantelung bezeichnet den Prozess,
in dem sich innerhalb des expandierenden
Apeirons eine unterscheidbare,
räumlich geschlossene Struktur bildet.

Das QUBLE besitzt eine annähernd
kugelförmige Gestalt.

\section*{Entstehung}

Die Expansion des Apeirons führt
zu lokalen Dichte- und
Geschwindigkeitsunterschieden.

Dynamische Scherung bewirkt
die Ausbildung einer Rotation.

Durch Rotationsschließung entsteht
eine räumlich begrenzte,
stabilisierte Struktur: das QUBLE.

\section*{Innere Ordnung}

Das QUBLE besitzt eine geschlossene
innere Apeiron-Strömung.

Seine räumliche Ordnung wird
durch Rotation und die fortlaufende
Expansion des Apeirons bestimmt.

Zwischen dem Inneren des QUBLES
und seiner Umgebung bestehen
Dichte- und Expansionsunterschiede.

\section*{Kontakt und Packungsbildung}

QUBLES bilden durch gegenseitige
Oberflächenberührung räumliche
Kontaktverbände.

Die Anordnung der Kontaktlagen
bestimmt die Geometrie der Packung.

Zwischen benachbarten QUBLES
entstehen Interstitien.

Die weitere räumliche Ordnung
führt zur Ausbildung der
Superdichten QUBLE-Packung (SDQP)
und der Dichten QUBLE-Packung (DQP).

\section*{Emergenzfolge}

\[
\begin{gathered}
\text{Apeiron}
\rightarrow
\text{Expansion}
\rightarrow
\text{Quantelung}
\\[0.4em]
\rightarrow
\text{QUBLE}
\rightarrow
\text{SDQP}
\rightarrow
\text{DQP}
\rightarrow
\text{Pfad}
\end{gathered}
\]

\vspace{0.9em}
\hrule
\vspace{0.45em}

{\footnotesize
MoE-Glossar:
\emph{Apeiron und Apeiron-Expansion},
\emph{SDQP},
\emph{DQP},
\emph{Interstitium},
\emph{Pfad}.
}

\end{document}
